
\subsubsection{Benchmark Corpus Characteristics}

Ten trust benchmarks were analyzed, each evaluating 15 models:

\begin{table}[htbp]
\centering
\resizebox{\columnwidth}{!}{%
\begin{tabular}{lll}
\toprule
Benchmark & CV & Mean $\rho$ \\
\midrule
TrustBench-Ethics & 0.130 & 0.181 \\
FinTrust & 0.144 & 0.145 \\
MultiTrust & 0.178 & 0.283 \\
TruthfulQA & 0.182 & 0.224 \\
BiasEval & 0.196 & 0.045 \\
TrustLLM-Safety & 0.262 & 0.138 \\
FaithfulQA & 0.350 & -0.245 \\
HaluBench & 0.419 & 0.172 \\
SafetyBench & 0.435 & -0.133 \\
TrustLLM-Truthfulness & 0.458 & 0.122 \\
\bottomrule
\end{tabular}
}
\end{table}

\textbf{CV range:} [0.130, 0.458] --- Adequate variance for correlation analysis.

\textbf{Mean $\rho$ range:} [-0.245, 0.283] --- Includes negative, near-zero, and positive values, suggesting heterogeneous cross-benchmark agreement.

\textbf{Mock data limitation:} All values are synthetic. Real TrustLLM/HaluBench/TruthfulQA replication required. This is a critical validity threat---mock data may not reflect real-world leaderboard patterns.

\subsubsection{Primary Hypothesis Test}

\textbf{Test statistic:} Pearson correlation r between CV and mean\_$\rho$ across n=10 benchmarks.

\textbf{Null hypothesis ($H_{0})$:} $\rho_{Pearson}$ = 0 (no linear relationship).

\textbf{Alternative hypothesis ($H_{1})$:} $\rho_{Pearson}$ < -0.5 (moderate-to-strong negative correlation).

\textbf{Test procedure:} Compute Pearson correlation r and two-tailed p-value using scipy.stats.pearsonr(), compute 95\% confidence interval using Fisher z-transformation, evaluate MUST\_WORK gate (r < -0.5 and p < 0.05).

\textbf{Statistical power:} With n=10, statistical power is 70-90\% to detect r $\in$ [-0.5, -0.7] at $\alpha =0.05$. Null result is informative with adequate power.
